\documentclass[troisieme]{fiche}
\metadonnees{20}{Un secret}{Bloc D — Rapporter et supposer}

\begin{document}
\entete

\objectifs{%
\textbf{Langue} : bilan de l'année ; thème.\par
\textbf{Texte} : Oscar Wilde, \og The Canterville Ghost \fg{} (1887), dernière page.\par
\textbf{Durée} : 1 h — exercices 20 min, texte et questions 25 min, version 15 min.}

%% ===================================================================
\exercice[12 min]{Bilan de l'année}

\rappel{\textbf{Rappel.} Chaque phrase porte sur un point vu cette année ; l'indication entre
parenthèses dit lequel.}

\consigne{Complétez.}

\begin{anglais}
\begin{enumerate}[label=\arabic*.,itemsep=2.5mm,leftmargin=*]
  \item Look! The ghost \trou{is walking} in the corridor again. \emph{(walk)}
  \item He \trou{has haunted} the house since 1584. \emph{(haunt)}
  \item Yesterday the twins \trou{threw} a pillow at him. \emph{(throw)}
  \item You \trou{mustn't} laugh at a ghost. \emph{(interdiction)}
  \item He was \trou{the unhappiest} ghost in England. \emph{(unhappy)}
  \item There wasn't \trou{any} oil left for his chains. \emph{(quantité)}
  \item The girl \trou{who} listened to him was called Virginia. \emph{(relatif)}
  \item The chains \trou{were heard} in the corridor every night. \emph{(hear, au passif)}
  \item She told them that she \trou{could not} say what had happened. \emph{(can, discours
    indirect)}
  \item If anyone \trou{asked} him, he \trou{would answer} that he only wanted to sleep.
    \emph{(ask, answer)}
\end{enumerate}
\end{anglais}

\note[La fiche d'où vient chaque phrase]{1 et 2, fiches 1 et 4 ; 3, fiche 3 ; 4, fiche 8 ;
5, fiche 12 ; 6, fiche 13 ; 7, fiche 14 ; 8, fiche 16 ; 9, fiche 17 ; 10, fiche 19.}

\note[Deux détails qui coûtent des points]{Phrase 8 : au passif, le temps ne change pas —
\emph{were heard} et non \emph{are heard}. Phrase 9 : après un verbe introducteur au passé,
\emph{can} devient \emph{could}.}

%% ===================================================================
\exercice[8 min]{Thème}
\consigne{Traduisez en anglais.}

\begin{quote}\itshape
Quand les Otis achetèrent Canterville Chase, on leur dit que la maison était hantée. Mr Otis
répondit qu'il ne croyait pas aux fantômes. Le premier soir, sa femme découvrit une tache de
sang près de la cheminée, et son fils la fit disparaître avec un détachant américain ; mais le
lendemain matin la tache revint. Le fantôme était plus malheureux qu'effrayant. Si quelqu'un
lui demandait ce qu'il voulait, il répondrait qu'il voulait seulement dormir. Ce fut Virginia, la
plus jeune des enfants, qui le comprit enfin.
\end{quote}

\lignes{10}

\corr{\textbf{Proposition.} When the Otises bought Canterville Chase, they were told that the
house was haunted. Mr Otis answered that he did not believe in ghosts. On the first evening
his wife found a stain of blood near the fireplace, and his son removed it with an American
stain remover; but the next morning the stain came back. The ghost was more unhappy than
frightening. If anyone asked him what he wanted, he would answer that he only wanted to sleep.
It was Virginia, the youngest of the children, who understood him at last.}

\note[Cinq points de contrôle]{\og on leur dit \fg{} : passif, \emph{they were told}.
\og répondit qu'il ne croyait pas \fg{} : discours indirect, le présent recule d'un cran.
\og plus malheureux qu'effrayant \fg{} : comparatif d'un adjectif long, \emph{more unhappy
than}. \og Si quelqu'un lui demandait\ldots il répondrait \fg{} : irréel, prétérit +
\emph{would}. \og Ce fut Virginia qui \fg{} : \emph{it was Virginia who}, la seule tournure
possible, l'anglais n'ayant pas de \og c'est\ldots que \fg{} aussi souple.}

\note[Le nom propre au pluriel]{\og les Otis \fg{} se dit \emph{the Otises} : l'anglais
accorde les noms de famille au pluriel, à la différence du français.}

%% ===================================================================
\newpage
\section{Texte}

\noindent\small\textit{Virginia Otis, quinze ans, a passé une nuit enfermée avec le fantôme de
Sir Simon, qu'elle a aidé à trouver enfin le repos. Elle n'a jamais dit à personne ce qui
s'était passé. Quatre ans plus tard, elle épouse le jeune duc de Cheshire.}\normalsize

\begin{texte}
The Duke and Duchess, after the \gl[a honeymoon]{honeymoon}{une lune de miel} was over, went
down to Canterville Chase, and on the day after their arrival they walked over in the afternoon
to the lonely \gl[a churchyard]{churchyard}{un cimetière} by the pine-woods. There had been a
great deal of difficulty at first about the \gl[an inscription]{inscription}{une inscription}
on Sir Simon's \gl[a tombstone]{tombstone}{une pierre tombale}, but finally it had been decided
to \gl[to engrave]{engrave}{graver} on it simply the initials of the old gentleman's name, and
the verse from the library window. The Duchess had brought with her some lovely roses, which
she \gl[to strew, strewed]{strewed}{joncher, répandre} upon the
\gl[a grave]{grave}{une tombe}, and after they had stood by it for some time they
\gl[to stroll]{strolled}{flâner} into the \gl[ruined]{ruined}{en ruine}
\gl[a chancel]{chancel}{un chœur (d'église)} of the old \gl[an abbey]{abbey}{une abbaye}. There
the Duchess sat down on a fallen \gl[a pillar]{pillar}{une colonne}, while her husband lay at
her feet smoking a cigarette and looking up at her beautiful eyes. Suddenly he threw his
cigarette away, took hold of her hand, and said to her, ``Virginia, a wife should have no
secrets from her husband.''

``Dear Cecil! I have no secrets from you.''

``Yes, you have,'' he answered, smiling, ``you have never told me what happened to you when
you were locked up with the ghost.''

``I have never told any one, Cecil,'' said Virginia, gravely.

``I know that, but you might tell me.''

``Please don't ask me, Cecil, I cannot tell you. Poor Sir Simon! I \gl[to owe]{owe}{devoir
(quelque chose à quelqu'un)} him a great deal. Yes, don't laugh, Cecil, I really do. He made
me see what Life is, and what Death \gl[to signify]{signifies}{signifier}, and why Love is
stronger than both.''

The Duke rose and kissed his wife lovingly.

``You can have your secret as long as I have your heart,'' he \gl[to murmur]{murmured}{murmurer}.

``You have always had that, Cecil.''

``And you will tell our children some day, won't you?''

Virginia \gl[to blush]{blushed}{rougir}.
\end{texte}
\source{Oscar Wilde, \og The Canterville Ghost \fg, dans \emph{Lord Arthur Savile's Crime and
Other Stories}, Londres, Osgood, 1891, ch.~\textsc{vii}.}

\partiel[15 min]{Questions}
\consigne{Répondez aux deux premières questions en français, aux deux dernières en anglais.}

\begin{enumerate}[label=\arabic*.,itemsep=2.5mm,leftmargin=*]
  \item Qu'a-t-on fini par faire graver sur la pierre tombale, et pourquoi si peu ?
    \lignes{2}
    \corr{Seulement les initiales du vieux monsieur et le vers inscrit sur la fenêtre de la
    bibliothèque. On ne s'est pas mis d'accord sur une inscription : rien de ce qu'on aurait pu
    écrire ne convenait à un fantôme de trois siècles.}
  \item Que demande Cecil à Virginia, et que répond-elle ?
    \lignes{3}
    \corr{Il lui demande ce qui lui est arrivé la nuit où elle a été enfermée avec le fantôme.
    Elle refuse de le dire, comme elle l'a toujours refusé, mais elle ajoute qu'elle doit
    beaucoup à Sir Simon : il lui a fait voir ce qu'est la vie, ce que signifie la mort, et
    pourquoi l'amour est plus fort que l'une et l'autre.}
  \item \en{The last line of the story is two words: ``Virginia blushed.'' Why does Wilde end
    there instead of telling us the secret?}
    \lignes{3}
    \corr{Because the secret is what has held the whole story together, and saying it would
    make it ordinary. The blush answers the question without a word: she will tell the
    children, one day. Wilde gives the reader exactly as much as the Duke gets, and no more.}
  \item \en{Find in the text one passive form, one comparative, and one sentence of reported
    speech.}
    \lignes{3}
    \corr{Passif : \emph{it had been decided to engrave on it simply the initials}, ou
    \emph{when you were locked up with the ghost}. Comparatif : \emph{why Love is stronger
    than both}. Discours indirect : \emph{you have never told me what happened to you}, où la
    question \og que t'est-il arrivé ? \fg{} est rapportée sans inversion.}
\end{enumerate}

\note[Fin de l'année]{Les trois formes de la question 4 ont été vues aux fiches 16, 11 et 17.
Ce sont elles qui portent la dernière page du conte : un passif pour l'inscription que
personne n'assume, un comparatif pour la seule chose que le fantôme ait apprise à Virginia,
une question rapportée pour ce qu'elle ne dira pas.}

%% ===================================================================
\section{Version}
\consigne{Traduisez le premier paragraphe du texte, de \emph{The Duke and Duchess} à
\emph{her beautiful eyes}.}

\lignes{16}

\corr{\textbf{Proposition de traduction.} Le duc et la duchesse, une fois la lune de miel
terminée, descendirent à Canterville Chase, et le lendemain de leur arrivée ils allèrent à
pied, dans l'après-midi, jusqu'au cimetière solitaire, près des bois de pins. Il y avait eu au
début beaucoup de difficultés au sujet de l'inscription à porter sur la pierre tombale de Sir
Simon, mais on avait fini par décider d'y graver simplement les initiales du vieux monsieur et
le vers de la fenêtre de la bibliothèque. La duchesse avait apporté de belles roses, qu'elle
répandit sur la tombe ; et, après être restés quelque temps auprès d'elle, ils gagnèrent à pas
lents le chœur en ruine de la vieille abbaye. Là, la duchesse s'assit sur une colonne
abattue, tandis que son mari s'étendait à ses pieds, fumant une cigarette et levant les yeux
vers son beau regard.}

\note[Choix de traduction 1]{\emph{after the honeymoon was over} : \emph{to be over} =
\og être fini \fg. La proposition entière se ramène en français à un participe — \og une fois
la lune de miel terminée \fg.}

\note[Choix de traduction 2]{\emph{it had been decided to engrave} : passif impersonnel, sans
agent nommé, parce que personne ne veut endosser la décision. Le français dit \og on avait
fini par décider \fg.}

\note[Choix de traduction 3]{\emph{looking up at her beautiful eyes} : \emph{to look up at}
dit d'un seul mouvement le regard qui se lève. Le français a besoin de deux termes — \og en
levant les yeux vers \fg.}

%% ===================================================================
\partiel{Lexique à mémoriser}
\begin{multicols}{2}\small
\begin{itemize}[label=--,itemsep=0.8mm,leftmargin=*]
  \item \textbf{a honeymoon} — une lune de miel
  \item \textbf{a churchyard} — un cimetière
  \item \textbf{a grave} — une tombe
  \item \textbf{to engrave} — graver
  \item \textbf{a husband} — un mari
  \item \textbf{a wife} — une épouse \textit{(pl. wives)}
  \item \textbf{a secret} — un secret
  \item \textbf{to owe} — devoir (de l'argent, un service)
  \item \textbf{to lock up} — enfermer à clef
  \item \textbf{to blush} — rougir
  \item \textbf{lovely} — charmant, ravissant
  \item \textbf{a ghost} — un fantôme \textit{(fiche 11)}
\end{itemize}
\end{multicols}

\end{document}
